\section{Parallel Execution Framework}

To maximize assembly throughput and align with Industry 4.0 efficiency standards, the system is designed to execute pick-and-place operations concurrently when the robotic arms are operating in mutually exclusive workspaces. Managing multiple heterogeneous manipulators requires a decoupled, asynchronous architecture to prevent control bottlenecks and ensure real-time responsiveness. This framework is built upon a distributed ROS 2 Action Server topology, orchestrated by a centralized state machine.

\subsection{Asynchronous Supervisory Control}
The core of the execution framework is the \texttt{UnifiedAssemblySupervisor} node. Implemented using Python's \texttt{asyncio} library and executing within a ROS 2 \texttt{MultiThreadedExecutor}, this node maintains the global state of the assembly sequence without blocking critical telemetry pipelines. 

When a sequence is received from the planning module, the Supervisor categorizes tasks based on their assigned manipulator (AR4 or ABB) and operation type (Independent Pick/Place or Handover). For parallel operations, the system bypasses standard sequential execution and spawns independent asynchronous worker threads (\texttt{execute\_ar4\_worker} and \texttt{execute\_abb\_worker}). 

By utilizing \texttt{ReentrantCallbackGroup} structures, these workers operate concurrently. They independently handle camera perception queries, grasp pose transformations, and hardware-specific coordinate offsets. This ensures that if one arm experiences a slight delay during perception or motion planning, the other arm continues its execution uninterrupted, maximizing the operational uptime of the cell.

\begin{figure}[htbp]
    \centering
    % PLACEHOLDER FOR IMAGE: Architecture diagram showing the UnifiedAssemblySupervisor dispatching asynchronous goals to the AR4 and ABB action servers.
    \rule{0.85\textwidth}{6.5cm} 
    \caption{Asynchronous dispatch architecture of the Unified Assembly Supervisor.}
    \label{fig:async_supervisor_architecture}
\end{figure}

\subsection{Distributed Action Server Architecture}
To decouple high-level decision-making from computationally intensive kinematic calculations, the system employs dedicated C++ Action Servers (\texttt{ar4\_task\_server} and \texttt{abb\_task\_server}). The Supervisor acts solely as an Action Client, sending parameterized \texttt{ExecuteTask} goals (e.g., "PICK", "PLACE", "HOME") and waiting for asynchronous feedback.

This distributed architecture offers two primary advantages:
\begin{enumerate}
    \item \textbf{Computational Isolation:} Inverse Kinematics (IK) solving and collision checking are offloaded to the individual C++ nodes, preventing the Python-based Supervisor from freezing during complex path computations.
    \item \textbf{Preemption and Cancellation:} The ROS 2 Action protocol allows the Supervisor to asynchronously cancel active goals. If the Zone Manager detects a proximity breach, the Supervisor instantly broadcasts a cancel request, halting the hardware controllers mid-trajectory before transferring control to the conflict resolution module.
\end{enumerate}

\subsection{Heterogeneous Trajectory Execution Constraints}
A major challenge in parallelizing heterogeneous robots is accounting for their differing mechanical reliabilities. The execution framework dictates specific motion constraints tailored to the capabilities of each manipulator.

The ABB IRB120, an industrial-grade manipulator, exhibits high repeatability and precision. Consequently, the \texttt{abb\_task\_server} utilizes standard joint-space interpolation (\texttt{move\_group->plan()}) for the entirety of its transit, allowing for rapid, fluid, and energy-efficient movements between waypoints.

Conversely, the AR4 is a research-grade arm, which inherently introduces minor positional uncertainties. To mitigate the risk of mechanical deviations—especially when interacting with dense brick clusters—the \texttt{ar4\_task\_server} enforces strict Cartesian path planning (\texttt{computeCartesianPath}). During the final approach and retreat phases of a pick or place, the AR4 end-effector is mathematically constrained to move in a perfectly vertical, straight line. If the planner cannot find a linear path for at least $95\%$ of the required distance, the trajectory is aborted. This strict constraint prevents the AR4 from executing unpredictable joint-space arcs that could inadvertently collide with adjacent workpieces.

\begin{figure}[htbp]
    \centering
    % PLACEHOLDER FOR IMAGE: Comparison diagram showing Joint-Space interpolation (curved path) vs. Cartesian path planning (straight vertical line) during a pick operation.
    \rule{0.7\textwidth}{5cm} 
    \caption{Trajectory generation constraints: Joint-space interpolation vs. Cartesian linear constraints.}
    \label{fig:trajectory_constraints}
\end{figure}